\chapter{Phương pháp đề xuất}
\label{chap:method}

\section{Tổng quan kiến trúc}
Hệ thống gồm năm thành phần nối tiếp (Hình~\ref{fig:arch}). Ba agent chuyên môn phân tích độc lập ba nguồn dữ liệu và
mỗi agent trả về một nhận định có cấu trúc. Conflict Analyzer đo mức mâu thuẫn giữa ba nhận định; nếu mâu thuẫn dưới
ngưỡng, các nhận định được chuyển thẳng tới Mediator; nếu vượt ngưỡng, Conflict Analyzer phân loại mâu thuẫn, phân tích
nguyên nhân gốc và kích hoạt tranh biện. Mediator tổng hợp các nhận định cuối cùng thành tín hiệu BUY, SELL hoặc
NEUTRAL, sau đó tín hiệu được quy đổi thành hành động giao dịch.

\begin{figure}[H]
\centering
\begin{tikzpicture}[
  node distance=0.55cm and 1.0cm,
  box/.style={draw,rounded corners=3pt,align=center,minimum height=1.0cm,text width=2.8cm,font=\small,fill=#1},
  box/.default=white,
  arr/.style={-{Stealth[length=2.5mm]},thick}]
  \node[box=blue!8] (mkt) {\textbf{Market Agent}\\giá, EMA, RSI,\\tín hiệu kỹ thuật};
  \node[box=blue!8,above=of mkt] (fin) {\textbf{Financial Agent}\\6 thống kê on-chain};
  \node[box=blue!8,below=of mkt] (sen) {\textbf{Sentiment Agent}\\tin tức ngày trước};
  \node[box=orange!12,right=1.0cm of mkt,text width=3.6cm] (ca) {\textbf{Conflict Analyzer}\\[2pt]\footnotesize$C = \alpha\overline{D}_{KL} + (1{-}\alpha)\mathrm{Var}$};
  \node[diamond,draw,aspect=2.2,right=0.9cm of ca,font=\small,inner sep=1pt,fill=orange!5] (dec) {$C \geq \tau$?};
  \node[box=red!8,above=1.0cm of dec,text width=3.3cm] (deb) {\textbf{RCA + Debate}\\2 vòng, trích dẫn bằng chứng,\\giảm độ tự tin};
  \node[box=green!10,below=1.0cm of dec,text width=3.2cm] (med) {\textbf{Mediator}\\[2pt]\footnotesize$\Sfinal=\sum_i d_i s_i \omega_i$};
  \node[box=green!25,right=0.8cm of med,text width=2.2cm] (out) {\textbf{BUY / SELL /}\\\textbf{NEUTRAL}};
  \draw[arr] (fin.east) -- (ca.west);
  \draw[arr] (mkt.east) -- (ca.west);
  \draw[arr] (sen.east) -- (ca.west);
  \draw[arr] (ca) -- (dec);
  \draw[arr] (dec) -- node[right,font=\small]{có} (deb);
  \draw[arr] (dec) -- node[right,font=\small]{không} (med);
  \draw[arr] (deb.east) -- ++(0.4,0) |- ([yshift=9pt]med.east);
  \draw[arr] (med) -- (out);
\end{tikzpicture}
\caption{Kiến trúc hệ thống đề xuất}
\label{fig:arch}
\end{figure}

\section{Các agent chuyên môn}
\label{sec:specialists}
Mỗi agent chuyên môn nhận một bản tóm tắt dữ liệu của riêng mình, trong đó \emph{mỗi mục bằng chứng nằm trên một dòng
riêng} (một giá trị chỉ số, một câu tin tức) để các bước sau có thể trích dẫn chính xác từng dòng. Trong thiết lập
thực nghiệm với dữ liệu CryptoTrade, với quyết định tại thời điểm mở cửa ngày $D$:
\begin{itemize}
  \item \textbf{Market Agent} đọc giá mở cửa 7 ngày gần nhất, EMA20, EMA50, RSI14 tính trên giá mở cửa tới ngày $D$, và
  ba tín hiệu kỹ thuật mà CryptoTrade dùng (giao cắt đường trung bình, MACD, Bollinger).
  \item \textbf{Financial Agent} đọc 6 thống kê giao dịch on-chain của ngày $D-1$ (số địa chỉ duy nhất, số giao dịch,
  tổng giá trị chuyển, phí trung bình, tổng dung lượng khối, số giao dịch coinbase) kèm phần trăm thay đổi so với ngày
  $D-2$.
  \item \textbf{Sentiment Agent} đọc các bài báo của ngày $D-1$, mỗi câu một dòng.
\end{itemize}

Mỗi agent trả về một nhận định gồm: tín hiệu $\in\{\text{BUY},\text{SELL},\text{NEUTRAL}\}$; vectơ niềm tin gồm hướng
$d_i\in\{+1,-1,0\}$ và cường độ $s_i\in[0,1]$; chuỗi lập luận; các đoạn bằng chứng; và siêu dữ liệu gồm độ nhiễu
(entropy), độ trùng lặp và thời điểm của dữ liệu. Độ nhiễu không do LLM tự chấm mà được tính từ dữ liệu:
\begin{itemize}
  \item Market: $e = 1 - |\text{RSI}_{14} - 50|/50$, tức RSI càng gần 50 thì tín hiệu kỹ thuật càng mơ hồ.
  \item Financial: $e = \min\!\big(\tfrac{1}{30}\cdot\text{trung bình}_k|\Delta_k\%|,\,1\big)$, tức các chỉ số on-chain biến
  động càng mạnh thì càng nhiễu.
  \item Sentiment: $e = 0{,}5$ (trung tính), vì dữ liệu CryptoTrade không có chỉ số Fear \& Greed.
\end{itemize}

\section{Conflict Analyzer}
\label{sec:conflict}
\paragraph{Ánh xạ niềm tin thành phân phối.} Gọi $x_i = d_i\cdot s_i\in[-1,1]$ là niềm tin có dấu của agent $i$.
Niềm tin được ánh xạ tuyến tính thành phân phối trên hai trạng thái giảm/tăng:
\begin{equation}
  P_i = \big(P_i^{\text{giảm}},\,P_i^{\text{tăng}}\big) = \big(0{,}5 - 0{,}5\,x_i,\;0{,}5 + 0{,}5\,x_i\big).
\end{equation}
Ánh xạ này liên tục theo $x_i$: một agent NEUTRAL ứng với phân phối đều, một agent BUY với độ tự tin thấp chỉ lệch nhẹ
khỏi phân phối đều.

\paragraph{Chỉ số mâu thuẫn.} Phân kỳ KL có làm mượt $\epsilon=10^{-9}$ được tính trên mọi cặp có thứ tự của $N=3$
agent, kết hợp với phương sai của các $x_i$:
\begin{align}
  \overline{D}_{KL} &= \frac{1}{N(N-1)}\sum_{i\neq j}\sum_{k}P_i^{k}\log\frac{P_i^{k}}{P_j^{k}+\epsilon},\\
  C &= \alpha\,\overline{D}_{KL} + (1-\alpha)\,\mathrm{Var}(x_1,\dots,x_N),
\end{align}
với $\alpha=0{,}6$. Khi $C\geq\tau=0{,}4$, hệ thống coi là có mâu thuẫn.

\paragraph{Phân loại mâu thuẫn.} Với $C\geq0{,}15$, các luật sau được áp dụng: \emph{Signal Conflict} khi đồng thời có
BUY và SELL; \emph{Temporal Conflict} khi chênh lệch trọng số thời gian giữa các agent vượt 0,4; \emph{Reliability
Conflict} khi chênh lệch độ nhiễu vượt 0,5; \emph{Redundancy Conflict} khi độ trùng lặp của một agent vượt 0,6. Với
$C<0{,}15$ hệ thống phân loại thẳng là không có mâu thuẫn. Kết quả phân loại được đưa vào bước phân tích nguyên
nhân gốc.

\paragraph{Phân tích nguyên nhân gốc (RCA).} LLM nhận các nhận định, loại mâu thuẫn và toàn bộ bằng chứng, rồi viết
1--3 nhận định giải thích vì sao các agent bất đồng. Mỗi nhận định phải được gắn nhãn dữ kiện, suy luận hoặc giả thuyết
và trích dẫn dòng bằng chứng cụ thể; quy trình kiểm tra nguồn được mô tả ở Mục~\ref{sec:grounding}.

\section{Tranh biện}
\label{sec:debate}
\paragraph{Cường độ phản biện.} Ở mỗi vòng, với mỗi agent $i$, hệ thống tính cường độ phản biện từ các agent còn lại:
\begin{equation}
  R_i = \mathrm{clip}\!\left(\frac{1}{1{,}5}\cdot\frac{1}{N-1}\sum_{j\neq i}\Big(w_b\,D_{\cos}(x_i,x_j) +
  w_e\,D_{J}(E_i,E_j) + w_l\,D_{L}(L_i,L_j)\Big),\;0,\;1\right),
\end{equation}
trong đó $D_{\cos}$ là khoảng cách cosine giữa niềm tin có dấu (bằng 0 khi cùng hướng, 2 khi ngược hướng), $D_J$ là
khoảng cách Jaccard giữa tập từ trong bằng chứng của hai agent, $D_L$ là khoảng cách Levenshtein chuẩn hóa giữa hai
chuỗi lập luận, và $w_b=w_e=w_l=1/3$.

\paragraph{Cập nhật độ tự tin.} Độ tự tin của agent giảm theo hàm mũ của cường độ phản biện:
\begin{equation}
  s_i^{\text{mới}} = s_i \cdot e^{-\beta R_i},\qquad \beta = 0{,}35.
\end{equation}
Phản biện yếu hầu như không ảnh hưởng, phản biện mạnh nhất làm độ tự tin giảm khoảng 30\% mỗi vòng. Cập nhật chỉ thay
đổi \emph{cường độ} $s_i$, không thay đổi \emph{hướng} $d_i$: tranh biện có thể làm một agent bớt chắc chắn nhưng không
làm agent đổi phe. Hệ quả của lựa chọn thiết kế này được phân tích ở Mục~\ref{sec:debate-analysis}.

\paragraph{Lập luận lại.} Mỗi agent nhận chuỗi lập luận cũ, lịch sử tranh biện tích lũy qua các vòng trước (theo cách
tổng hợp của MADAM-RAG~\cite{wang2025madamrag}), bằng chứng của mọi agent và cường độ phản biện, rồi viết lại lập luận.
Lời nhắc ràng buộc mỗi agent giữ đúng vai trò chuyên môn để tránh các agent hòa thành một giọng sau vài vòng. Mặc định
có 2 vòng.

\section{Kiểm tra nguồn trích dẫn}
\label{sec:grounding}
Mọi nhận định do RCA và tranh biện sinh ra đi qua hai lớp kiểm tra:
\begin{enumerate}
  \item \textbf{Kiểm tra tất định bằng mã}: đúng định dạng, mã bằng chứng và mã dòng tồn tại; nội dung trích dẫn được
  lấy nguyên văn từ dòng được chỉ định, không để mô hình tự chép lại số liệu.
  \item \textbf{Kiểm tra ngữ nghĩa bằng một lời gọi LLM riêng}: đối chiếu từng nhận định với toàn bộ bằng chứng, trả về
  ``được hỗ trợ'', ``không được hỗ trợ'' hoặc ``mâu thuẫn''. Nếu kiểm tra lỗi hoặc không có kết quả, nhận định bị từ chối.
\end{enumerate}
Mô hình được tạo lại tối đa một lần khi bị từ chối. Nếu một cập nhật tranh biện vẫn bị từ chối, \emph{toàn bộ} cập
nhật của agent ở vòng đó bị hủy, kể cả độ tự tin, để một lập luận không có căn cứ không ảnh hưởng tới quyết định. Ngày
đó vẫn được giao dịch theo trạng thái đã hoàn tác và được gắn cờ để phân tích riêng.

\section{Mediator}
\label{sec:mediator}
Trọng số của agent $i$ là trọng số cơ bản nhân với ba hệ số phạt:
\begin{equation}
  \omega_i = b_i\,(1 - e_i)\,(1 - r_i)\,e^{-\gamma\Delta t_i},
\end{equation}
với $b_i=1$ là trọng số cơ bản, $e_i$ là độ nhiễu, $r_i$ là độ trùng lặp (0,05 cho Market; 0,1 cho Financial và
Sentiment), $\Delta t_i$ là tuổi của dữ liệu tính bằng giây so với thời điểm quyết định và $\gamma=10^{-5}$. Tín hiệu tổng
hợp và quyết định cuối là:
\begin{equation}
  \Sfinal = \sum_{i=1}^{N} d_i\,s_i\,\omega_i,\qquad
  \text{tín hiệu} = \begin{cases}
    \text{BUY} & \text{nếu } \Sfinal > 0{,}05,\\
    \text{SELL} & \text{nếu } \Sfinal < -0{,}05,\\
    \text{NEUTRAL} & \text{ngược lại.}
  \end{cases}
\end{equation}

\paragraph{Nhận xét về trọng số thực tế.} Dữ liệu on-chain và tin tức là của ngày $D-1$, nên $\Delta t = 86\,400$ giây
và hệ số thời gian là $e^{-0{,}864}\approx 0{,}42$; dữ liệu giá có $\Delta t = 0$. Vì vậy, với cùng độ nhiễu, trọng số của
Financial và Sentiment chỉ bằng khoảng 40\% trọng số của Market; riêng Sentiment với $e=0{,}5$ có trọng số cố định
khoảng $0{,}5\times0{,}9\times0{,}42\approx0{,}19$. Market Agent do đó chiếm ưu thế về cấu trúc trong $\Sfinal$. Đây là
một giả định thiết kế ngầm sẽ được xem xét lại trong hướng phát triển (Chương~\ref{chap:conclusion}).

\section{Quy đổi tín hiệu thành hành động}
Tín hiệu được quy đổi thành hành động trong môi trường giao dịch của CryptoTrade theo cách bài báo quy đổi các chiến
lược luật: BUY $\mapsto +0{,}5$ (dùng 50\% tiền mặt để mua), SELL $\mapsto -0{,}5$ (bán 50\% lượng BTC đang giữ),
NEUTRAL $\mapsto 0$. Nếu một agent chuyên môn lỗi, hệ thống giữ nguyên vị thế trong ngày đó.

\section{Kiểm soát rò rỉ thông tin}
Quyết định tại giá mở cửa ngày $D$ chỉ dùng giá tới giá mở cửa ngày $D$ và dữ liệu on-chain, tin tức của ngày $D-1$,
trùng khớp với thông tin mà agent CryptoTrade nhận được. Một bài kiểm thử tự động xác nhận sự trùng khớp này trên toàn
bộ 61 ngày của giai đoạn tăng giá. Mỗi lời nhắc cũng yêu cầu mô hình chỉ dùng bằng chứng được cung cấp, không dùng
kiến thức về diễn biến sau thời điểm dự báo.
